\vspace{-2mm}
\section{Introduction}
\label{sec:intro}
3D scene editing is a long-standing problem in computer vision. By offering intuitive control over content creation, it plays a crucial role in diverse applications such as visual effects~\cite{zhang2025advancing3dgaussiansplatting}, augmented and virtual reality (AR/VR)~\cite{madhavaram2024trainingfreeapproach3d}, game design~\cite{xu2024sketch2sceneautomaticgenerationinteractive}, and digital twins~\cite{melnik2025digitaltwingenerationvisual}.
%
Recent advancements in 3D scene representation and reconstruction~\cite{nerf, 3dgs} have sparked a surge of interest in editing 3D scenes using text prompts. Existing methods~\cite{in2n, gaussctrl, lee2025editsplatmultiviewfusionattentionguided} typically combine neural reconstruction with text-driven image editing in a \textit{reconstruct-edit-refit} pipeline. Given a collection of posed images and a textual instruction, these approaches first \textit{reconstruct} the scene using a neural 3D representation, then apply 2D generative \textit{editing} to the rendered images, and finally \textit{refit} (or re-optimize) the 3D representation to incorporate the edits.
%
While this pipeline achieves impressive visual quality and alignment with instruction prompts, it is limited by slow per-scene optimization and retraining, which limits practical usability. Moreover, these methods often introduce multi-view inconsistencies, revealing a substantial gap between current 3D editing frameworks and the need for fast, interactive scene manipulation in real-world applications.

Conversely, the rapid development of Large Reconstruction Models (LRMs)~\cite{lrm, gslrm, noposplat, mvsnerf, pixelsplat, dust3r} offers a more efficient paradigm for 3D scene reconstruction by amortizing computation across large datasets and removing per-scene optimization. Notably, Gaussian-based LRMs enable rapid, photorealistic rendering suited for real-time applications, thanks to the Gaussian splatting technique.

Motivated by these advancements, we propose {\ourmethod}, a feed-forward approach that reconstructs and edits a 3D scene in a single pass.
%
Beginning with a pair of unposed views, we edit them using a 2D image editor. These potentially inconsistent edited views serve as inputs for {\ourmethod}, which directly predicts 3D Gaussian splats aligned with the edited views while eliminates the inconsistency, seamlessly unifying reconstruction and editing without the need for per-scene optimization.
% Our method takes as input a few unposed images of a scene, where a reference view is edited according to the textual instruction. From these potentially inconsistent views, \ourmethod directly predicts the 3D Gaussian splats representing the edited scene

Training such a feed-forward reconstruction-editing model introduces two practical challenges. First, ground-truth edited images for supervision are unavailable. To address this, we leverage SAM2-based recoloring~\cite{ravi2024sam} for cross-view-consistent supervision, providing reliable labels and masks. We find that our model, though trained only on a recoloring task, it effectively handles images edited by 2D diffusion models at inference. Second, edited inputs can remain semantically and visually inconsistent across views. To tackle this, we design an asymmetric input scheme that pairs recolored reference views with other views in their original appearance, encouraging the network to fuse and align edited and unedited observations.

Moreover, the lack of established benchmarks limits fair and scalable evaluation in 3D scene editing. We introduce {\ourbench}, a standardized benchmark built on the DL3DV~\cite{ling2023dl3dv10klargescalescenedataset} test split to systematically assess multi-view consistency and inference efficiency, setting a foundation for reproducible and rigorous comparisons.
%
Our contributions are summarized as follows:
\begin{itemize}
\item We introduce {\ourmethod}, a feed-forward framework that simultaneously reconstructs 3D scenes and generates edited Gaussian splats aligned to a text prompt.
\item We develop a SAM2-based recoloring technique for view-consistent supervision during training and introduce an asymmetric input scheme to bolster robustness against cross-view inconsistencies.
\item We present {\ourbench}, a scene-level, multi-view 3D editing benchmark built on the DL3DV test split.
\item Extensive quantitative and qualitative experiments show {\ourmethod} is significantly faster than previous approaches, achieving higher visual quality and semantic consistency while demonstrating practical real-world potential.
\end{itemize}




% %-------------------------------------------------------------------------
% %
% Scene editing is a long-standing problem in computer vision, enabling intuitive control over content creation across VFX, AR/VR authoring, game design, and digital-twin applications~\cite{}.
% %
% Recent advancements in 3D reconstruction techniques (\eg, Neural Radiance Fields (NeRFs) and 3D Gaussian Splatting), together with generative models that make instruction-guided editing practical from simple text prompts~\cite{}, is driving increased interest in 3D scene editing .
% %
% Several research efforts~\cite{in2n, gaussctrl} propose methods that combine neural reconstruction with text-driven image editing to enable 3D scene editing.
% Given a collection of posed images and a textual editing instruction, these approaches first reconstruct the original scene with a neural 3D representation, then apply 2D generative editing to the rendered images, and finally re-optimize or retrain the 3D representation to incorporate the edits. 
% %
% Although they deliver high visual quality and renderings aligned with the input prompts, this reconstruct–edit–refit pipeline faces practical limitations. In particular, slow optimization and retraining hinder usability, and multi-view consistency often remains fragile, leaving a gap between current 3D editing capabilities and real-world needs for rapid, interactive scene manipulation.

% % LRM介绍
% The advent of Large Reconstruction Models (LRMs)~\cite{lrm, gslrm, noposplat, mvsnerf, pixelsplat, dust3r} provides a significantly more efficient route to 3D scene reconstruction by amortizing computation across large datasets and eliminating expensive per-scene optimization.
% %
% Among them, Gaussian-based LRMs directly predicts Gaussian primitives, positions, covariances, colors, and opacities, from sparse posed views, enabling rapid, photorealistic rendering that suits real-time and practical applications.
% %
% However, current LRMs primarily aim for faithful reconstruction of static scenes rather than flexiblemanipulation or editing.
% %
% Applying LRMs on scene editing is nontrivial, as training LRMs requires multi-view consistent supervision (e.g., edited ground truth), which is rarely available; producing perfectly matched edits across views is impractical, leaving the acquisition of reliable labels as a persistent bottleneck.

% In this paper, we propose {\ourmethod}, a feed-forward approach that reconstructs and edits a 3D scene in a single pass under textual instructions.
% %
% Given a pair of unposed, sparse, instruction-edited views with remaining cross-view inconsistencies, {\ourmethod} directly predicts instruction-aligned 3D Gaussian splats in a single pass, eliminating per-scene optimization and the separation between reconstruction and editing.

% Training and evaluating such a feed-forward reconstruction–editing model raises three practical challenges. (1) Ground-truth edited images for non-reference views are unavailable, and cross-view inconsistencies persist regardless of the chosen 2D editor; we therefore derive cross-view–consistent supervision via SAM2~\cite{}-based recoloring, which provides reliable labels/masks. (2) During inference, the edited inputs remain semantically and visually inconsistent across views. To mitigate this, we adopt an asymmetric input scheme during training to manually create inconsistency: the recolored reference view is paired with other views kept in their original appearance, encouraging the network to learn to fuse and align edited and unedited observations. (3) Fair evaluation is hindered by small, ad-hoc benchmarks. We thus automate a standardized benchmark on DL3DV~\cite{}, termed {\ourbench}, to systematically assess multi-view consistency and inference efficiency, enabling reproducible comparisons.

% Our contributions are summarized as follows:
% \begin{itemize}
%     \item We propose {\ourmethod}, a feed-forward method that reconstructs the 3D scene and generates edited Gaussian splats based on the input instruction.
%     \item We propose a SAM2-recolor-based and asymmetric input training strategy to obtain view-consistent supervision during training.
%     \item We introduce {\ourbench}, a scene-level, multi-view 3D editing benchmark built on the DL3DV test split for scene editing.
%     \item Results show that our method accomplish much faster than existing apporachs with better visual quality and semantic consistency, highlighting its potential for real-world application.
% \end{itemize}

% \XT{